% Generated draft: requires mathematical and source review.
\documentclass[11pt]{article}
\usepackage[margin=25mm]{geometry}
\usepackage{fontspec}
\setmainfont{Latin Modern Roman}
\usepackage{amsmath,amssymb,mathtools,unicode-math}
\setmathfont{Latin Modern Math}
\usepackage{enumitem,needspace,xurl,hyperref}
\setlength{\parindent}{0pt}
\setlength{\parskip}{5pt}
\setlength{\emergencystretch}{3em}
\newcommand{\N}{\mathbb{N}}
\newcommand{\Z}{\mathbb{Z}}
\newcommand{\Q}{\mathbb{Q}}
\newcommand{\R}{\mathbb{R}}
\newcommand{\C}{\mathbb{C}}
\newcommand{\F}{\mathbb{F}}
\newcommand{\A}{\mathbb{A}}
\newcommand{\PP}{\mathbb{P}}
\newcommand{\E}{\mathbb{E}}
\newcommand{\eps}{\varepsilon}
\newcommand{\dd}{\,\mathrm{d}}
\begin{document}
\section*{UnsolvedMath ID: 600007}
\textbf{Source problem code:} catalogue entry
\subsection{Definitions and mathematical statement}
Let $f\in\R[x,y]$ have $0$ as a regular value, meaning $\nabla f(x,y)\ne0$ at every real point with $f(x,y)=0$. Let $\gamma$ be an oval component of this zero set, interpreted as a compact smooth simple closed component. The source assumes that positive nearby levels foliate an outer neighborhood: for some $\delta>0$, the relevant curves at $0<\varepsilon<\delta$ are disjoint leaves filling an exterior annular neighborhood of $\gamma$.

For each tangent line $\ell$ at $p\in\gamma$, the local paired intersections of these leaves with $\ell$ are assumed to be exchanged by one projective involution of $\ell$, independent of $\varepsilon$. In an affine coordinate on $\ell$, such an involution is locally represented by a nonidentity fractional linear map $t\mapsto(at+b)/(ct+d)$ with $ad-bc\ne0$ whose square is the identity. It exchanges the two intersections for every sufficiently small applicable level and fixes their limiting contact point $p$. This states the natural paired-intersection meaning; if more local branches or tangencies are allowed, the exact source pairing needs clarification.

The requested conclusion is that $\gamma$ is an ellipse and the nearby curves are conics in one pencil. A pencil is a family $uQ_0+vQ_1=0$, $[u:v]\in\mathbb{RP}^1$, for two fixed independent homogeneous quadratic forms in projective plane coordinates. The source writes $\gamma_\varepsilon=\{f=\varepsilon\}$, but a polynomial can have other level components away from the chosen oval. It must be clarified whether the conclusion concerns only the outer leaves or the entire real/algebraic level sets; these are not interchangeable. Likewise ``oval'' here does not automatically insert an unstated strict-convexity hypothesis needed in some outer-billiard formulations.

The stored research summary claims a stronger description: after choosing an elliptic quadratic function $Q$ with $\gamma=\{Q=c_0\}$, one has $f=P\circ Q$ for a one-variable polynomial $P$, and the relevant leaves are the pencil $\{Q=c\}$, generated projectively by the ellipse and the double line at infinity. This is a stored claim, not a theorem independently proved or reviewed in this draft. In particular $P\circ Q$ can have several level components, so the stronger formula does not justify calling an entire arbitrary level set a single conic without resolving the leaf convention.

The summary describes three proposed proof steps: polynomial invariance along a tangent line forces a projective involution fixing its contact point, written $t\mapsto-t/(1-\lambda t)$, to have $\lambda=0$; the resulting central symmetries give outer-billiard invariance and invoke Tabachnikov's 2008 Theorem 1; after an affine circle normalization, the twist map is written $(\theta,r)\mapsto(\theta+2\arccos(1/r),r)$, and a rotation-number argument plus polynomiality yields $P(x^2+y^2)$. The hypotheses of the cited theorem and global leaf conclusions are not reproduced, so these steps remain an attributed proof claim requiring checking.

The registry labels the record solved and its research field says \texttt{SOLVED-BY-YOU}; that label does not mean the present assistant supplied a proof or that human mathematical review occurred. The original 2015 source presents an open problem, and the provenance notes accessed 29 July 2026 mark current status and rights as \texttt{NEEDS\_REVIEW}. These source limitations are retained while the statement-only draft proceeds.
\subsection{Short English statement}
Does a polynomial outer foliation whose tangent-line intersection pairings are projective involutions force an ellipse and a pencil of conics? The stored solution claim uses outer billiards, but its review status, convexity assumptions and local-leaf versus full-level interpretation require checking.
\subsection{Sources}
\begin{itemize}
\item Local UnsolvedMath registry, permanent ID $600007$, exact source code \texttt{catalogue entry}, ``Baker's Dozen---A converse Desargues theorem,'' including its stored research summary and provenance/release-review notes.
\item Tabachnikov, \emph{A Baker's Dozen of Problems} (2015), Problem 3, AMR-indexed source list, \url{https://amj.math.stonybrook.edu/html-articles/Files-2015-2024/14-01/}, accessed 29 July 2026 according to the record; extracted from publicly accessible publisher HTML.
\item The stored research summary invokes Tabachnikov, Theorem 1, Pacific Journal of Mathematics \textbf{235} (2008), 89--92, without a supplied article title, URL, or full hypotheses. It is cited only as part of that stored proof claim.
\end{itemize}
\end{document}
